\section{Magnetic, superconducting and Dirac targets}\label{app:E.20}

Clean aluminum needs no first-principles calculation for its loss function near the pair-breaking edge, only three numbers: the density of states $\mathcal{N}_{\mathrm{F}}$ (free-electron value), the Fermi velocity $v_{\mathrm{F}} = \SI{2.0e6}{m.s^{-1}}$ and the gap $\Delta_{\mathrm{BCS}} = \SI{0.18}{meV}$; its conserving response is a two-dimensional integral that runs in minutes. Amorphous WSi, at the other extreme, needs molecular dynamics, DFT, the RPA and a Bogoliubov--de Gennes solution on each quenched cell. The magnets and Dirac materials lie in between. Their rates are written in the book through a few model parameters (exchange constants, anisotropies, form factors, Dirac velocities and background screening), and the calculations below supply those parameters and the responses that go beyond the model. \Cref{tab:E20-params} lists the calculations with their costs.

\subsection{The conserving response of aluminum}

The conserving density response of \cref{sec:15.1} adds a phase-mode term to the bare quasiparticle bubble, and the two cancel at $\vect{q} = 0$. We want its vertex fraction $\mathfrak{v}$ on the grid of that section: $y = \omega/\Delta_{\mathrm{BCS}}$ from 2 to 50 and $qv_{\mathrm{F}}/\Delta_{\mathrm{BCS}}$ from its kinematic minimum $y\,v_{\mathrm{F}}/v_{\mathrm{max}}$ to $10^3$, together with the full $\Imag\varepsilon^{\mathrm{cons}}$, its Kramers--Kronig transform and the loss function $L^{\mathrm{cons}}$ in its resummed form at the edge. The electrons of a typical film are diffusive on the scale of the probe (\cref{sec:15.3}), so the calculation includes impurity scattering. On the truncation label $\mathfrak{T}$ (\cref{sec:8.5}) it is $\mathcal{X}$ = RPA in the Coulomb channel with the pairing ladder, on the BCS reference, quasiclassical (constant $\mathcal{N}_{\mathrm{F}}$, spherical Fermi surface, $T = 0$), clean or with Born impurity scattering.

Each clean bubble is an integral over the angle $\mu$ between $\vect{k}$ and $\vect{q}$ and over the band energy $\xi$, with $\xi' = \xi + v_{\mathrm{F}}q\mu + q^2/2m_e$, $E = (\xi^2 + \Delta_{\mathrm{BCS}}^2)^{1/2}$, and the pair-creation coherence factors $\frac12[1 - (\xi\xi' - \Delta_{\mathrm{BCS}}^2)/(EE')]$ for the density vertex and $\frac12[1 + (\xi\xi' + \Delta_{\mathrm{BCS}}^2)/(EE')]$ for the phase vertex. Gauss--Legendre quadrature in $\mu$ and adaptive quadrature in $\xi$ evaluate the bubble. The imaginary parts are taken at the roots of $E + E' = \omega$, solved analytically, with no broadening, and the real parts are principal values with the pole subtracted. The phase--phase bubble diverges logarithmically in the $\xi$ cutoff, but only its difference from the value at $(0,0)$ enters, so the cutoff is taken at $10^3\Delta_{\mathrm{BCS}}$ and the logarithm is cancelled analytically; the value at $(0,0)$ itself must be computed with the same cutoff as the gap equation, where it equals the inverse pairing coupling, or the Goldstone cancellation fails and $\mathfrak{v}$ acquires a spurious constant. The Kramers--Kronig integral runs on a logarithmic grid to $10^3qv_{\mathrm{F}}$, with the Lindhard normal state above $qv_{\mathrm{F}}$. The sign of $\ii0$ on the antiresonant pole is that of the retarded function, with which the time-ordered form agrees only for $\omega > 0$ at $T = 0$, and the bare and conserving bubbles are never mixed between $\Imag\varepsilon$ and $\Real\varepsilon$, since mixing them breaks positivity.

Impurities enter through the Born self-energy and the impurity ladder of \cref{sec:15.3}, whose normal state is the Mermin function \cref{eq:4.5-mermin} at the relaxation rate $\gamma_{\mathrm{rel}} = v_{\mathrm{F}}/\mathfrak{l}_e$. The table is computed in the clean limit and for $\mathfrak{l}_e = \SI{1}{\micro m}$, \SI{300}{nm} and \SI{30}{nm}, with the edge values divided by $\mathfrak{g}_{\mathrm{I}}$ times the Mermin normal state, as in \cref{tab:15-edge}.

The conserving response must vanish at $q \to 0$ for every $\omega$, and with the ladder switched off it must return the bare bubble with the density coherence factor~\cite[Eq.~(7)]{gan2026mean}. $\mathfrak{v}$ must be stable to $10^{-3}$ under doubling of both grids, the Kramers--Kronig tail must stay below $10^{-3}$ of $\Real\varepsilon$, and the loss-function table must be converged to \SI{1}{\percent}. With impurities the normalized edge ratio has to reproduce the conserving values of \cref{sec:15.3} at the same mean free paths, and $L$ must still vanish at $2\Delta_{\mathrm{BCS}}$. The table replaces the exclusion of the cells nearest the edge in \cref{sec:15.2}, for the mean free path of the film at hand.

\subsection{Amorphous WSi}

A W$_x$Si$_{1-x}$ film has $k_{\mathrm{F}}\mathfrak{l}_e$ of a few (1.8--5.5 from its measured transport), and at $q = 1$ to \SI{5}{keV}, where $q$ is comparable to both $k_{\mathrm{F}}$ and the screening wavenumber, no relaxation-time form of $\varepsilon$ is valid (\cref{sec:15.3}). The rate needs the configurational average $\langle L\rangle(\vect{q},\omega)$ of the loss function, not the loss function of the averaged $\varepsilon$, for $q = 0.5$ to \SI{5}{keV} and $\omega$ from $2\Delta_{\mathrm{BCS}}$ to \SI{50}{eV}, together with the distribution $P_\Delta$ of the local gap. On $\mathfrak{T}$ the normal state is $n_B = 0$, $\mathcal{X}$ = RPA with local fields, averaged over amorphous configurations, and the gap is a Bogoliubov--de Gennes mean field with an attractive interaction. The superconducting factor is applied locally with $P_\Delta$; Anderson's theorem justifies this local treatment where the mean free path is far shorter than the coherence length.

A machine-learned potential (\cref{app:E.11}) is trained on PBE molecular dynamics of 216-atom cells at the film composition, $x \approx 0.2$--0.25, and 512- to 1000-atom cells are melt-quenched at no more than \SI{e13}{K.s^{-1}}, at least five times, and validated on the measured density and the W--Si and W--W pair distributions. Each quench gets a DFT ground state (SG15, with the tungsten 5s5p semicore in the valence) at $\mathrm{\Gamma}$ plus a $2^3$ grid, and a non-self-consistent run with at least three times as many bands as are occupied, so that the semicore states, which contribute above about \SI{35}{eV}, are in the response. The Kubo--Greenwood conductivity of each quench, $\sigma_1(\omega)$, is the optical control, compared with film ellipsometry and sheet resistance. The RPA dielectric matrix (Yambo or BerkeleyGW, \cref{app:E.6}) is computed directly at the reciprocal vectors of the supercell, which are $2\pi/L \approx \SI{0.5}{keV}$ apart for $L = \SI{25}{\angstrom}$, so no folding is needed above \SI{0.5}{keV}. The $\vect{G}$ cutoff is converged to $q + \SI{3}{keV}$ and the bands to $\omega_{\mathrm{max}} + \SI{30}{eV}$. Local fields are large in a disordered cell, and the loss function is $\Imag[-(\varepsilon^{-1})_{\vect{0}\vect{0}}]$, never $\Imag[-1/\varepsilon_{\vect{0}\vect{0}}]$. A Wannier tight-binding model of each cell, solved in the Bogoliubov--de Gennes approximation with the attractive interaction fitted to the measured film $T_c$, gives $\Delta(\vect{r})$ and $P_\Delta$. The $f$-sum per $q$ must hold to \SI{2}{\percent} and the measured $\sigma_1$ must be reproduced; the error bar is the spread of $\langle L\rangle$ across quenches and from 512 to 1000 atoms, which must stay below \SI{5}{\percent}. The RPA matrices, at tens of gigabytes each for 1000 atoms, dominate the cost (\cref{tab:E20-params}), and the result enters \cref{sec:15.3}.

\subsection{Holstein absorption with Eliashberg functions}

Above $2\Delta_{\mathrm{BCS}}$ a superconductor absorbs a dark photon by breaking a pair and emitting a phonon, with the conductivity $\sigma_1^{\mathrm{S}}(\omega)$ of \cref{eq:15.4-holstein}. The rate within a few meV of the edge depends on the shape of $\alpha^2F_{\mathrm{tr}}(\omega')$ below a few meV, where the free-electron model holds $1.25\times10^4$ times the weight of the $\omega'^4$ law (\cref{sec:15.4}). We therefore want $\sigma_1^{\mathrm{S}}(\omega)/\sigma_n$ for aluminum and niobium with $\alpha^2F_{\mathrm{tr}}$ from first principles, also in its strong-coupling form, where the pair weight becomes $\int\dd E\,[N^{\mathrm{El}}(E)N^{\mathrm{El}}(E') + P^{\mathrm{El}}(E)P^{\mathrm{El}}(E')]$ with the Eliashberg density of states and pair amplitude of \cref{sec:15.4-strong}. On $\mathfrak{T}$ it is $n_g = 2$, one electron--phonon vertex and one current vertex, with isotropic Eliashberg dressing.

EPW~\cite{ponce2016epw} computes the vertices by DFPT on an $8^3$ grid and interpolates them to fine grids of $60^3$ $\vect{k}$ and $30^3$ $\vect{Q}$ for aluminum and $40^3$ and $20^3$ for niobium, for both $\alpha^2F$ and the transport function with its weight $1 - \cos\theta$. The $30^3$ grid does not resolve the $\omega'^4$ regime, so the transport function below a few meV is recomputed on a refined grid of small $\vect{Q}$ around $\mathrm{\Gamma}$, and it must follow the $\omega'^4$ law there. That law holds only if the vertex obeys the acoustic sum rule \cref{eq:6.2-asr}, which makes it vanish for $\vect{Q}\to0$ between states on the Fermi surface, so the sum rule is checked on the interpolated vertex. The isotropic Eliashberg equations~\cite{margine2013anisotropic} are solved on the real axis, since Pad\'e continuation produces artifacts near $2\Delta_{\mathrm{BCS}}$, and \cref{eq:15.4-holstein} is then integrated with a phonon step no larger than $\Delta_{\mathrm{BCS}}/20$. The transport coupling $\lambda_{\mathrm{tr}}$ must match the value 0.39 from the resistivity of aluminum, against which the free-electron model gives 0.32 (\cref{sec:15.4}); $T_c$ must come within \SI{10}{\percent} of experiment with a single Coulomb pseudopotential $\mu^*$, $\sigma_1$ must be stable to \SI{5}{\percent} under doubling of the fine grids, and the normal-state $\sigma_1$ must lie within \SI{20}{\percent} of the measured optical conductivity of aluminum between 0.2 and \SI{3}{eV}. \Cref{sec:15.4-strong} leaves open which renormalization factors survive in the strong-coupling form; in the normal metal the factor is $(1 + \lambda_{\mathrm{tr}})^{-2} = 0.52$ in the rate, and the calculation reports the superconducting result with and without the factor $1/(1 + \lambda_{\mathrm{tr}})$ in the current. The difference between the two results then measures the effect of this factor.

\subsection{Magnet inputs from DFT~\texorpdfstring{$+\,U$}{+U}}\label{app:E.20-magnet}

The magnon rates of \cref{sec:14.1,sec:14.2} need exchange constants $J^{\mathrm{ex}}_{ij}$, single-ion anisotropies and $g_{\mathrm{ion}}$ tensors. For YIG, NiO, MnO, Cr$_2$O$_3$ and NiPS$_3$ these come from inelastic neutron or resonance fits where they exist, and otherwise from DFT $+\,U$ (\cref{app:E.9}), in the Heisenberg convention of the book, $J^{\mathrm{ex}} > 0$ antiferromagnetic with each pair counted once. The Hubbard $U$ is computed by linear response~\cite{timrov2018hubbard}, and the exchange constants by four-state energy mapping or by the Green-function method of TB2J~\cite{he2021tb2j} on Wannier functions, in the magnetic unit cell at the measured lattice constant, shell by shell until a further shell no longer moves the computed magnon band by more than its spread over $U$. The spread over $U$ is the main systematic, and the measured ordering temperature and, where available, the measured magnon band are the checks. The orbital moments of Cr$^{3+}$ and Ni$^{2+}$ enter $g_{\mathrm{ion}}$ and are computed with spin--orbit coupling. Spin--orbit coupling also gives the Dzyaloshinskii--Moriya constants of Cr$_2$O$_3$, by four-state mapping with the two spins of each bond placed in the plane normal to the component sought. The spin model of \cref{sec:14.5} needs the Dzyaloshinskii--Moriya constants, together with the exchange and anisotropy constants, for $T \ll T_{\mathrm{CM}}$. For a target without ultrasonic data, the magnetoelastic constants $B_1$ and $B_2$ of \cref{sec:14.2} follow from the strain derivative of the magnetic anisotropy energy, computed with spin--orbit coupling at a few strained cells, with strains of $\pm\SI{0.5}{\percent}$ as for the deformation potentials; where ultrasonic data exist, they check the computed constants.

The same linear spin-wave code that serves the two-magnon response below gives the uniform modes that absorption into magnons needs (\cref{sec:14.3}). For a sample of finite size the uniform modes carry the demagnetizing field of its shape, and in a body that is not an ellipsoid the dipolar lattice sums over the real sample spread the weight over several magnetostatic modes, so the code takes the sample shape and the dipolar sums as inputs. In the ferromagnetic limit the uniform mode must lie exactly at the Larmor frequency, and in the two-sublattice antiferromagnet the uniform transverse weight per antiparallel pair must equal $|\vect{S}|\omega_{\mathrm{AFMR}}/(2\omega_E + \omega_A)$. YIG, with 80 atoms in its magnetic cell, dominates the cost.

\subsection{Two-magnon responses}

At linear spin-wave order a ferromagnet at $T = 0$ has no two-magnon density response, while an antiferromagnet has a two-magnon continuum whose weight relative to the one-magnon line is 0.146 to 0.029 for $|\vect{S}| = 1/2$ to 5/2 in a simple-cubic lattice (\cref{sec:14.1}). The quartic terms of the Holstein--Primakoff expansion add a two-magnon interaction of relative order $1/|\vect{S}|$, which \cref{sec:14.2} resums as a ladder in the pair channel of total momentum $\vect{q}$; its poles below the continuum are two-magnon bound states. Two results are open: the existence of a two-magnon bound state in NiO, and the interacting response $S^{(2m)}(\vect{q},\omega)$ in \si{eV^2}. We want the interacting response, and the bound-state energies in meV with their weights, for YIG, NiO, MnF$_2$ and, as an uncontrolled benchmark, $\alpha$-RuCl$_3$. On $\mathfrak{T}$ the reference is linear spin-wave theory with the pair ladder, and for noncollinear order the cubic vertex is added.

Linear spin-wave theory with Colpa's diagonalization~\cite{colpa1978diagonalization} (SpinW~\cite{toth2015linear} or Sunny), with the inputs of \cref{app:E.20-magnet}, gives the reference. The quartic vertex is derived symbolically from the spin Hamiltonian, and the Holstein--Primakoff and Dyson--Maleev vertices agree only on shell, so one of them is used throughout. The ladder is one linear solve per $(\vect{q},\omega)$ on $30^3$ to $40^3$ grids, about $3\times10^7$ operations per $\vect{q}$, and exact diagonalization of the two-magnon sector in small cells checks it. The zero-point reduction $\Delta S$, a sum over Bogoliubov coefficients, is converged to $10^{-3}$ on a $40^3$ grid. The sum rule on $\langle\vect{S}^2\rangle$ must hold to \SI{1}{\percent}, and the bound-state energies must be stable to \SI{1}{\percent} when the grid is doubled. YIG, with 20 iron ions per primitive cell, needs the full multi-sublattice transformation, and a matrix that is not positive definite signals a wrong classical ground state. The results enter \cref{sec:14.2}.

\subsection{Orbital form factors}

Ions with unquenched orbital moments, the $4f$ ions and some $3d$ ions with strong spin--orbit coupling, scatter with form factors beyond the dipole form $\langle j_0\rangle + c_2\langle j_2\rangle$ at $q \gtrsim 1/\langle r\rangle_{4f} \approx \SI{3}{keV}$. Expanding $\ee^{\ii\vect{q}\cdot\vect{r}}$ in spherical Bessel functions and applying the Wigner--Eckart theorem inside the ground $J$ multiplet reduces the spin and orbital matrix elements to the radial integrals $\langle j_K(q)\rangle = \int r^2\dd r\,R_{n\ell}(r)^2j_K(qr)$ with $K$ even up to $2\ell = 6$, times angular coefficients. The integrals for $K = 0,\dots,6$ (dimensionless) and the multiplet-projected matrix elements are wanted for the ions of the magnetic targets of \cref{sec:14.1}, with the crystal field treated as a perturbation within $J$; crystal-field mixing of $J$ multiplets matters for the light rare earths. On $\mathfrak{T}$ this is a single-ion calculation in the ground multiplet at clamped nuclei, and the magnetic response that uses it is at linear spin-wave order. Dirac--Fock calculations with the General-purpose Relativistic Atomic Structure Package (GRASP) or fully relativistic DFT $+\,U$ Wannier projections give the radial functions. Dark matter that couples to the orbital moment couples to the orbital current, not to $\vect{\ell}$ times the density, and the operator is built from the orbital current. The integrals must satisfy $\langle j_0(0)\rangle = 1$ and the orbital-moment sum rules to $10^{-3}$, and agree with the tabulated dipole form factors to \SI{2}{\percent} for $q \le \SI{1}{keV}$.

\subsection{Finite-momentum optical activity}

In a chiral crystal the current response has an antisymmetric part odd in $\vect{q}$, $\ii\epsilon_{\alpha\beta\gamma}\mathcal{G}^\gamma_{\mathrm{OA}}(\vect{q},\omega)$ in $S_{j^\alpha j^\beta}$ (\cref{eq:14.5-GOA}), which changes sign between the two enantiomers and gives the difference of their rates in \cref{sec:14.5}. DFPT gives only its slope at $q = 0$, the natural optical activity~\cite{zabalo2023natural}. We want $\mathcal{G}_{\mathrm{OA}}(\vect{q},\omega)$, with the units of $S_{j^\alpha j^\beta}$, and the full current tensor at $\pm\vect{q}$ for tellurium or $\alpha$-quartz and for the magnetoelectric crystal Cr$_2$O$_3$, at $q \le \SI{10}{keV}$ and $\omega \le \SI{50}{eV}$. On $\mathfrak{T}$ the longitudinal part is at $\mathcal{X}$ = RPA with local fields and the transverse part at independent particles with local fields, both with spin--orbit coupling. Noncollinear DFT with fully relativistic SG15 pseudopotentials gives the states. The current vertex $\langle\vect{k}+\vect{q}|\{\vect{p} + \ii[V_{\mathrm{NL}},\vect{r}],\ee^{\ii\vect{q}\cdot\vect{r}}\}/2|\vect{k}\rangle$ needs the nonlocal commutator (\cref{app:E.7}) and is evaluated on Wannier-interpolated states in the same gauge at $+\vect{q}$ and $-\vect{q}$, and Yambo adds the local fields of the longitudinal channel. The sign of the odd part follows the sign convention for $\vect{q}$ of \cref{app:A}. The odd part must vanish to numerical noise in silicon, which is centrosymmetric, and its linear term must reproduce the DFPT optical activity; the grid goes from $16^3$ to $24^3$ until the odd part changes by less than \SI{5}{\percent}.

\subsection{The ZrTe\texorpdfstring{$_5$}{5} edge}

The excitonic binding of a gapped Dirac material is set by $\mathrm{Ry}_{\mathrm{D}} = \Delta_{\mathrm{D}}\tilde\alpha^2/4$, the Rydberg of a pair with reduced mass $\Delta_{\mathrm{D}}/(2v_{\mathrm{D}}^2)$ and coupling $\alpha/\varepsilon_{\mathrm{bg}}$, with $\tilde\alpha = \alpha/(\varepsilon_{\mathrm{bg}}v_{\mathrm{D}})$ the effective coupling of \cref{eq:7.3-alpha}. For ZrTe$_5$ the central parameters are those of the measured crystal structure (\cref{sec:16.1-zrte5}, the first ZrTe$_5$ row of \cref{tab:E20-dirac}): the half-gap $\Delta_{\mathrm{D}} = \SI{11.8}{meV}$, the principal velocities $1.1\times10^{-3}$, $4.5\times10^{-4}$ and $1.0\times10^{-3}$ with geometric mean $v_{\mathrm{D}} = 7.9\times10^{-4}$, the energy cutoff $\tilde\Lambda = \SI{0.191}{eV}$, $g_{\mathrm{deg}} = 2$, and $\varepsilon_{\mathrm{bg}} = 50.5$, the harmonic mean of the dielectric tensor. Then $\tilde\alpha \approx 0.18$, $\mathrm{Ry}_{\mathrm{D}} \approx \SI{0.10}{meV}$, and the exciton radius is $a_{\mathrm{D}} = 2v_{\mathrm{D}}/(\tilde\alpha\Delta_{\mathrm{D}}) \approx \SI{1.4e3}{\angstrom}$. A $24^3$ grid over the Brillouin zone has a spacing of \SI{1.9e-2}{\angstrom^{-1}}, 27 times $1/a_{\mathrm{D}}$, so an ab initio BSE cannot resolve the threshold. Instead we want $L(\vect{q},\omega)$ of the anisotropic massive-Dirac model for $|\vect{q}| \le \Lambda_{\mathrm{D}}$ and $2\Delta_{\mathrm{D}} \le \omega \le 2\Delta_{\mathrm{D}} + \SI{50}{meV}$ (at $\vect{q} \to 0$; at finite $\vect{q}$ the invariant pair energy replaces $\omega$), with the bound-state energies and weights. This $\omega$ range contains the region up to 1.33 times the edge where the Sommerfeld factor exceeds 1.5 (\cref{sec:16.1}).

$G_0W_0$ with spin--orbit coupling on the 12-atom cell ($12\times12\times4$ grid, refined near $\mathrm{\Gamma}$) fits the direction-resolved velocities $v_{\mathrm{D},i}$, the half-gap, the cutoffs $\Lambda_{\mathrm{D},i}$ and the background tensor $\varepsilon_{\mathrm{bg},ii}$. Since $\mathrm{Ry}_{\mathrm{D}}$ lies far below the optical phonons, the right background is the static tensor with the ionic part, and the $G_0W_0$ stage computes it that way as its own check on the tabulated tensor; remote bands belong in $\varepsilon_{\mathrm{bg}}$, not also in the $GW$ screening. The threshold region, within about 10 $\mathrm{Ry}_{\mathrm{D}}$ of $2\Delta_{\mathrm{D}}$, is the Elliott--Sommerfeld problem of an anisotropic hydrogen atom with masses $\Delta_{\mathrm{D}}/(2v_{\mathrm{D},i}^2)$ and dielectric tensor $\varepsilon_{\mathrm{bg},ii}$, solved on a radial grid. Away from threshold the BSE of the model, with the Coulomb interaction screened by $\varepsilon_{\mathrm{bg}}$ plus the RPA Dirac polarization, runs on a logarithmically graded radial grid from $k_{\mathrm{min}} \le 0.1/a_{\mathrm{D}}$ to $\Lambda_{\mathrm{D}}$, about 200 radial points times 590 Lebedev directions, matched to the analytic solution over 5 to 10 $\mathrm{Ry}_{\mathrm{D}}$; a uniform grid would give a spurious absence of bound states. The spread of $v_{\mathrm{D},i}$ and $\Delta_{\mathrm{D}}$ between one-shot and eigenvalue-self-consistent $GW$, and between the measured-structure and relaxed rows, is propagated through both stages. On $\mathfrak{T}$ this is $\mathcal{X}$ = BSE with static screening, the controlled first correction in $\tilde\alpha$; beyond it no controlled method exists at $\tilde\alpha$ of order one, so that stage gives only an estimate of the theory uncertainty. $\mathrm{Ry}_{\mathrm{D}}$ and the first three bound states must be stable to \SI{5}{\percent} when $k_{\mathrm{min}}$ is halved and the radial points doubled, the Sommerfeld factor must be continuous to \SI{2}{\percent} across the matching window, and the model $f$-sum must hold to $10^{-3}$. The model stages take minutes and the $G_0W_0$ about $10^4$ core-hours; the same run gives $\varepsilon_{\mathrm{bg}}$ of Bi$_2$Se$_3$.

\begin{table}[t]
  \centering
  \caption{Dirac parameters of ZrTe$_5$~\cite[Tables~1 and 2]{geilhufe2020dirac}: half-gap, principal velocities (units of $c$), diagonal background dielectric tensor and degeneracy. The first row is the measured crystal structure, the central row of \cref{sec:16.1-zrte5}; the second is the relaxed cell.}
  \label{tab:E20-dirac}
  \small
  \begin{tabular}{lllll}
    \toprule
    material & $\Delta_{\mathrm{D}}$ (meV) & $v_{\mathrm{D},x}, v_{\mathrm{D},y}, v_{\mathrm{D},z}$ ($10^{-4}$) & $\varepsilon_{\mathrm{bg},xx}, \varepsilon_{\mathrm{bg},yy}, \varepsilon_{\mathrm{bg},zz}$ & $g_{\mathrm{deg}}$ \\
    \midrule
    ZrTe$_5$, measured structure & 11.8 & 11, 4.5, 10 & 308.4, 20.75, 126.1 & 2 \\
    ZrTe$_5$, relaxed cell & 15.6 & 11, 4.4, 9.1 & 308.4, 20.75, 126.1 & 2 \\
    \bottomrule
  \end{tabular}
\end{table}

\begin{table}[t]
  \centering
  \caption{Magnetic, superconducting and Dirac calculations: targets, main grids, acceptance tests and cost in core-hours. KK, Kramers--Kronig.}
  \label{tab:E20-params}
  \small
  \setlength{\tabcolsep}{4pt}
  \begin{tabular}{p{0.17\textwidth}p{0.17\textwidth}p{0.23\textwidth}p{0.21\textwidth}p{0.09\textwidth}}
    \toprule
    quantity & targets & grids and parameters & acceptance test & cost \\
    \midrule
    vertex fraction $\mathfrak{v}$, $L^{\mathrm{cons}}$ & Al, clean and $\mathfrak{l}_e = \SI{1}{\micro m}$, \SI{300}{nm}, \SI{30}{nm} & $y \in [2,50]$; $qv_{\mathrm{F}}/\Delta_{\mathrm{BCS}} \le 10^3$; KK to $10^3qv_{\mathrm{F}}$ & $q\to0$ cancellation; bare limit; edge values of \cref{sec:15.3} & minutes \\
    $\langle L\rangle$, $P_\Delta$ & a-W$_x$Si$_{1-x}$ & 512--1000 atoms, ${\ge}5$ quenches & $f$-sum \SI{2}{\percent}; film $\sigma_1$ & \num{1.6e4}--\num{1.6e5} \\
    $\sigma_1^{\mathrm{S}}(\omega)$, $\alpha^2F_{\mathrm{tr}}$ & Al, Nb & $\vect{k}$ $60^3$, $\vect{Q}$ $30^3$ (Al), refined small $\vect{Q}$ & $\lambda_{\mathrm{tr}}$ 0.39; acoustic sum rule; $T_c$ \SI{10}{\percent}; optical $\sigma_1$ & $10^3$ \\
    $J^{\mathrm{ex}}$, anisotropy, Dzyaloshinskii--Moriya, $B_1$, $B_2$, $g_{\mathrm{ion}}$ & YIG, NiO, MnO, Cr$_2$O$_3$, NiPS$_3$ & magnetic cell; strained cells & $T_{\mathrm{N}}$, $T_{\mathrm{C}}$; magnon bands; Larmor; two sublattices & $10^3$--$5\times10^3$ \\
    $S^{(2m)}$, bound states & YIG, NiO, MnF$_2$ & $30^3$--$40^3$ & $\langle\vect{S}^2\rangle$ \SI{1}{\percent}; grid doubling & $10^2$ \\
    $\langle j_K(q)\rangle$ & $4f$ and spin--orbit ions & $K \le 6$ & $\langle j_0(0)\rangle = 1$; dipole tables & ${<}10$ per ion \\
    $\mathcal{G}_{\mathrm{OA}}$ & Te or $\alpha$-SiO$_2$; Cr$_2$O$_3$ & $16^3$--$24^3$ & zero in Si & $10^3$--$10^4$ \\
    $L$ near $2\Delta_{\mathrm{D}}$ & ZrTe$_5$ & log radial, 590 directions & $\mathrm{Ry}_{\mathrm{D}}$ \SI{5}{\percent}; $f$-sum & $10^4$ \\
    \bottomrule
  \end{tabular}
\end{table}
